% ------- Conjuntos -------

\newcommand{\set}[1]{   % Descripción conjunto cualquiera (llaves)
    \left\{
        #1
    \right\}
}

    % Conjuntos de números

\newcommand{\N}{    % Números naturales
    \mathbb{N}
}
\newcommand{\Z}{    % Números enteros
    \mathbb{Z}
}
\newcommand{\Q}{    % Números racionales
    \mathbb{Q}
}
\newcommand{\I}{    % Números irracionales
    \mathbb{I}
}
\newcommand{\R}{    % Números reales
    \mathbb{R}
}
\newcommand{\C}{    % Números complejos
    \mathbb{C}
}
\newcommand{\F}{    % Campo arbitrario ?
    \mathbb{F}
}
\newcommand{\Sf}{	% Esfera
	\mathbb{S}
}

% Conjuntos vectoriales

\newcommand{\GL}{
	\mathrm{GL}
}

\newcommand{\SL}{
	\mathrm{SL}
}

\newcommand{\Ort}{
	\mathrm{O}
}

\newcommand{\SO}{
	\mathrm{SO}
}



    % Topología

\newcommand{\kleinb}{
	\mathbb{K}
}

\newcommand{\T}{
	\mathbb{T}
}

\newcommand{\teich}{
	\mathcal{M}
}

\newcommand{\simtimes}{\mathbin{%
	\stackrel{\sim}{\smash{\times}\rule{0pt}{0.9ex}}
	% Twisted product
}}

	% Mapping Class Groups

\newcommand{\Mod}{
	\mathrm{Mod}
}

\newcommand{\Homeo}{
	\mathrm{Homeo}
}

\newcommand{\Diff}{
	\mathrm{Diff}
}


    % Conjuntos asociados a funciones de Álgebra Lineal


\newcommand{\traza}[1]{
	\textbf{traza}
	\left(#1\right) 
}

% ------- Funciones -------

    % Normas

\providecommand{\abs}[1]{   % Valor absoluto
    \left\vert
        #1
    \right\vert
}
\providecommand{\norm}[1]{  % Norma
    \left\lVert
        #1
    \right\rVert
}
\providecommand{\prodint}[1]{   % Producto interior
    \left\langle
        #1
    \right\rangle
}

    % Probabilidad

\newcommand{\Esp}[1]{   % Esperanza matemática
    E
    \left[
        #1
    \right]
}
\newcommand{\Prob}[1]{  % Probabilidad
    P
    \left[
        #1
    \right]
}

    % dimensiones de Álgebra Lineal

\renewcommand{\dim}{   % Dimensión
    \mathrm{dim} \,
}

    % Números complejos
\newcommand{\repart}[1]{
    \mathrm{Re}
    \left(
        #1
    \right)
}
\newcommand{\impart}[1]{
    \mathrm{Im}
    \left(
        #1
    \right)
}

% ----- Letras especiales ------

\newcommand{\eps}{  % Épsilon bonita (siempre varépsilon)
    \varepsilon
}

\newcommand{\vphi}{
    \varphi
}

% ----- Operadores diferenciales ------

\newcommand{\Cinf}{
	\mathcal{C}^\infty
}
